\section{Related Work} \label{sect:related_work}

{\bf On-device AI.}
Various techniques facilitate the adoption of AI on edge devices. FACIL~\cite{facil} applies processing-in-memory to overcome the memory capacity and bandwidth limitations of edge devices. DecDEC~\cite{decdec} introduces an aggressive low-bit quantization scheme for efficient inference. These works accelerate LLM inference in general, and can be applied to \proposed for further speedup. AppAgent~\cite{appagent} mimics human-like interactions to function without system access. Mobile-Agent~\cite{mobile_agent, mobile_agent_v2} leverages visual capabilities to identify and locate elements in the devices. Overall, these target multi-modal LLM-based agents, whereas \proposed focuses on conventional agentic workflows with text-based LLMs.


{\bf KV cache reuse.}
A rich body of prior work makes use of pre-computed KV caches to reduce runtime latency. Prefix caching~\cite{vllm} reuses KV caches from previous sequences when there are exact prefix matches to the incoming request. However, its reuse is limited when there are mismatches early on in the prompt. Prompt Cache~\cite{prompt_cache} reuses the full KV cache even when there are token matches that do not start at the beginning of the input, and CacheBlend~\cite{cacheblend} selectively recomputes KV caches to maximize reuse while maintaining accuracy. \proposed reuses KV caches by reconstructing the prompts at the text level, targeting the semantic similarity, without any recomputation at runtime.

{\bf Speculative decoding.} \label{page:related_work_specdec}
To facilitate speculative decoding~\cite{sps, specdec}, works like Eagle~\cite{eagle} reduce the training cost of draft LLMs by reusing target model's logits, whereas self-speculative decoding~\cite{draft_and_verify} reuses portions of the target model as draft models, completely removing the need for any retraining. \proposed is distinct from these works in that the lookup table used as the draft model requires no training, and its lightweight, constant time lookup speed allows for speedup directly proportional to the draft token latency. PLD~\cite{pld} constructs an LUT from the user prompt to accelerate input-grounded tasks like summarization. \proposed is different from PLD in that it pinpoints specific parts of the prompt to construct an LUT, and applies selective decoding to avoid the multi-token tax when possible.
